\documentclass[aspectratio=169,10pt]{beamer}
\usepackage[UTF8,fontset=none]{ctex}
\usepackage{booktabs}
\usepackage{array}
\usetheme{tutorial-tech-light}

\title[VSP-Beamer Tech]{VSP-Beamer Tech\\电路风格幻灯片模板}
\subtitle{Title band with circuit traces / cards / vector crest}
\VSPsetspeaker[Speaker]{VSP-Beamer Team}{ShanghaiTech University \quad SIST}
\date{November 2026}
\VSPsetupfootline{VSP-Beamer Tech}{Template}{VSP-Beamer Team}

\begin{document}
\VSPtitleframe

\begin{frame}{Contents}
  \tableofcontents
\end{frame}

\section{Page Structure}

\begin{frame}[t]{Title band and footline}
\begin{columns}[T,onlytextwidth]
  \begin{column}{.56\textwidth}
    标题栏右侧是电路走线与校徽，下沿是蓝紫渐变细线；标题上方的小字是当前章节名。
    \begin{itemize}
      \item 西文使用 Latin Modern，中文使用 Noto Sans CJK SC，与其他主题统一
      \item \alert{alert} 为紫色强调，\textbf{textbf} 为标准加粗
      \item 页脚左侧是 \texttt{\textbackslash VSPsetupfootline} 的前两格，右侧是页码
    \end{itemize}
  \end{column}
  \begin{column}{.40\textwidth}
    \begin{vspquote}{使用方法}
      \texttt{\textbackslash usetheme\{tutorial-tech-dark\}}\\
      \texttt{\textbackslash usetheme\{tutorial-tech-light\}}

      深浅两版接口相同，文稿可以直接切换。
    \end{vspquote}
  \end{column}
\end{columns}
\end{frame}

\section{Cards and Numbers}

\begin{frame}[t]{Key numbers}
\begin{columns}[T,onlytextwidth]
  \begin{column}{.31\textwidth}\VSPstat{33\,ms}{Reconstruct}{per frame}\end{column}
  \begin{column}{.31\textwidth}\VSPstat[TechViolet]{100\,Mbit/s}{Transmit}{5G downlink}\end{column}
  \begin{column}{.31\textwidth}\VSPstat[TechGreen]{22.1\,ms}{Decode}{laptop GPU}\end{column}
\end{columns}

\vspace{6pt}
\texttt{\textbackslash VSPstat[颜色]\{数字\}\{标签\}\{备注\}} 把关键数字放在卡片上。
\end{frame}

\begin{frame}[t]{Callouts}
\begin{columns}[T,onlytextwidth]
  \begin{column}{.32\textwidth}
    \begin{vspquote}{Note}默认强调色。\end{vspquote}
    \begin{vspquotepurple}{Violet}旁注与延伸。\end{vspquotepurple}
  \end{column}
  \begin{column}{.32\textwidth}
    \begin{vspquoteblue}{Blue}方法或定义。\end{vspquoteblue}
    \begin{vspquotegreen}{Green}已完成的结果。\end{vspquotegreen}
  \end{column}
  \begin{column}{.32\textwidth}
    \begin{vspquotered}{Warning}需要注意的地方。\end{vspquotered}
    \begin{block}{Beamer block}标准 \code{block} 同样适用。\end{block}
  \end{column}
\end{columns}
\end{frame}

\section{Figures and Code}

\begin{frame}[t]{Figure on a card and a table}
\begin{columns}[T,onlytextwidth]
  \begin{column}{.40\textwidth}
    \centering
    \begin{VSPcard}
      \includegraphics[width=.86\linewidth]{\VSPasset{ShanghaiTech_Logo.pdf}}
    \end{VSPcard}
  \end{column}
  \begin{column}{.56\textwidth}
    \begin{Table}[Variants]
      \footnotesize
      \begin{tabular}{@{}lll@{}}
        \toprule
        Theme & Canvas & Logo \\
        \midrule
        \texttt{tutorial-tech-dark} & near-black blue & white \\
        \texttt{tutorial-tech-light} & light grey & red \\
        \bottomrule
      \end{tabular}
    \end{Table}
  \end{column}
\end{columns}
\end{frame}

\begin{frame}[t,fragile]{Code and formula}
\begin{columns}[T,onlytextwidth]
  \begin{column}{.48\textwidth}
\begin{lstlisting}[language=bash]
make list
make check
\end{lstlisting}
  \end{column}
  \begin{column}{.48\textwidth}
    \begin{vspquoteblue}{Loss}
      \[
        \mathcal{L} = \sum_{i=1}^{N} \bigl\lVert \hat{I}_i - I_i \bigr\rVert_1
          + \lambda\,\mathcal{L}_{\mathrm{SSIM}}
      \]
    \end{vspquoteblue}
  \end{column}
\end{columns}
\end{frame}

\VSPendframe{VSP-Beamer \quad Tech theme}
\end{document}
